%% OmniLaTeX Advanced Bibliography Gallery
%%
%% Every citation style and biblatex feature in one document.
%% Compile with:
%%   latexmk -xelatex main.tex

\documentclass[
    language=english,
    doctype=article,
    institution=none,
]{../../omnilatex}

\RequirePackage{config/document-settings}

%% ─────────────────────────────────────────────────────
%% BibLaTeX Configuration
%% ─────────────────────────────────────────────────────
%% backend=biber is the default; shown here for clarity.
%% style=authoryear provides author-year citations.
%% dashed=true replaces repeated author names with a dash.
%% sorting=nyt sorts by name, year, title.
\usepackage[
    backend=biber,
    style=authoryear,
    dashed=true,
    sorting=nyt,
    maxbibnames=99,
    minbibnames=3,
    uniquelist=false,
    doi=true,
    isbn=false,
    url=false,
    eprint=false,
]{biblatex}

%% Add the main bibliography and a local supplementary file.
\addbibresource{../../bib/bibliography.bib}
\addbibresource{local-refs.bib}

%% ─────────────────────────────────────────────────────
%% Inline bibliography commands
%% ─────────────────────────────────────────────────────
\DeclareCiteCommand{\footfullcitehelper}[\footnote]
  {\usebibmacro{prenote}}
  {\printnames{labelname}%
   \setunit{\addcomma\space}%
   \printfield[citetitle]{labeltitle}%
   \setunit{\addcomma\space}%
   \printfield{year}}
  {\multicitedelim}
  {\usebibmacro{postnote}}

\begin{document}

\maketitle

\tableofcontents

% =====================================================
\section{Citation Styles Overview}
\label{sec:citation-styles}
% =====================================================

BibLaTeX provides a rich set of citation commands, each formatting the
reference differently.  This section demonstrates every standard citation
command available in the \texttt{authoryear} style.

\subsection{Basic Citation Commands}

The standard \verb|\cite| command produces a compact citation suitable for
most contexts.  See \cite{einstein1905} for the original formulation of
special relativity.

Parenthetical citations place the reference in parentheses:
\parencite{knuthTeXbook1986}.

Textual citations integrate the author name into the running text.
\textcite{goossensLaTeXCompanion1993} remains the definitive guide to
\LaTeX{} typesetting.

Footnote citations place the full reference in a footnote:
\footcite{mckinneyPythonDataAnalysis2018}.

Smart citations adapt to the context---they produce parenthetical or
textual output depending on the surrounding syntax:
\smartcite{ramalhoFluentPython2015}.

The \verb|\autocite| command is style-dependent and respects the global
citation style setting: \autocite{birdNaturalLanguageProcessing2009}.

\subsection{Author and Year Extraction}

Sometimes you need only part of the citation.  Extract just the author
name: \citeauthor{diracPrinciplesQuantumMechanics1981}.  Extract just the
publication year: \citeyear{mollenhauerHandbuchDieselmotoren2007}.  Extract
just the title: \citetitle{baehrThermodynamikGrundlagenUnd2016}.

\subsection{Full Citations}

The \verb|\fullcite| command reproduces the complete bibliography entry
inline: \fullcite{einstein1905}.

% =====================================================
\section{Citation Qualifiers and Notes}
\label{sec:qualifiers}
% =====================================================

Citation commands accept optional prenote and postnote arguments to add
textual qualifiers.  The syntax is
\verb|\cite[<prenote>][<postnote>]{<key>}|.

See, for example, \cite[see][p.~42]{knuthTeXbook1986} for a discussion of
the typesetting algorithm.  For a specific page range:
\cite[pp.~891--921]{einstein1905}.  A single optional argument provides a
postnote only: \cite[p.~144]{PTBSI}.

Parenthetical qualifiers work identically:
\parencite[for an overview][Chapter~3]{goossensLaTeXCompanion1993}.

% =====================================================
\section{Multiple Work Citations}
\label{sec:multiple}
% =====================================================

When citing several works at once, use the plural forms of citation
commands.  The \verb|\cites| command handles multiple keys:
\cites{einstein1905}{knuthTeXbook1986}{goossensLaTeXCompanion1993}.

Parenthetical multiple citations:
\parencites{ramalhoFluentPython2015}{mckinneyPythonDataAnalysis2018}{birdNaturalLanguageProcessing2009}.

Footnote multiple citations:
\footcites{diracPrinciplesQuantumMechanics1981}{baehrThermodynamikGrundlagenUnd2016}.

Each plural command also accepts per-entry qualifiers:
\cites[see][p.~5]{einstein1905}[and][p.~100]{knuthTeXbook1986}.

% =====================================================
\section{Date Tracking}
\label{sec:date-tracking}
% =====================================================

BibLaTeX tracks dates precisely.  Extract just the year:
\citeyear{einstein1905}.  When an entry contains a full date field, you can
access it with \verb|\citedate|.  The Einstein paper was published on
\citedate{einstein1905}.

This is especially useful for works where the exact date matters, such as
standards: the ISO standard was published on
\citedate{internationalorganizationforstandardizationISO821720172017}.

% =====================================================
\section{Citation Counting}
\label{sec:citation-counting}
% =====================================================

BibLaTeX provides a \verb|\citecounter| that tracks how many distinct
works have been cited.  After the first citation, the counter reads
\citecounter.  This increments with each new unique entry cited throughout
the document and can be used in custom formatting or statistics.

% =====================================================
\section{Bibliography Sections and Categories}
\label{sec:categories}
% =====================================================

BibLaTeX allows the bibliography to be split into sections by type,
keyword, or other field.  This is essential for documents that
distinguish primary sources from secondary literature.

\subsection{Type-Based Filtering}

Only journal articles:
\printbibliography[type=article,heading=subbibliography,title={Journal Articles}]

Only books:
\printbibliography[type=book,heading=subbibliography,title={Books}]

\subsection{Keyword-Based Filtering}

Entries tagged with the \texttt{primary} keyword (primary sources):
\printbibliography[keyword=primary,heading=subbibliography,title={Primary Sources}]

Entries tagged with the \texttt{reference} keyword:
\printbibliography[keyword=reference,heading=subbibliography,title={Reference Works}]

\subsection{Sub-bibliography with TOC Integration}

A sub-bibliography that appears in the table of contents:
\printbibliography[heading=subbibintoc,title={Full Bibliography (integrated in TOC)}]

% =====================================================
\section{Inline and Footnote Bibliography}
\label{sec:inline}
% =====================================================

The \verb|\footfullcite| command places a complete bibliography entry in a
footnote at the point of citation.  This is useful when a full reference
must appear immediately\footfullcite{einstein1905}.

Another inline example: this technique is common in legal and historical
writing\footfullcite{diracPrinciplesQuantumMechanics1981}.

% =====================================================
\section{Custom Entry Types}
\label{sec:custom-types}
% =====================================================

BibLaTeX supports custom entry types beyond the standard @book, @article,
etc.  The following examples demonstrate @online, @software, and @dataset
entries defined in the supplementary bibliography.

\subsection{Online Sources}

Online references include URL and access-date metadata:
\cite{GBV}.

Another online entry: \cite{TUB}.

And a software documentation page: \cite{mathworksCreateSimpleClass2020}.

\subsection{Software}

Software citations are increasingly common in reproducible research.
Our analysis tool is available as \cite{omnilatex2026}.

A related software package: \cite{biber2024}.

\subsection{Datasets}

Scientific datasets deserve proper citation.  The benchmark dataset is
available at \cite{kaggleGraphBenchmark2025}.

A supplementary dataset: \cite{openalexCitationGraph2025}.

% =====================================================
\section{Related Entries}
\label{sec:related}
% =====================================================

BibLaTeX's \texttt{related} field links entries that are editions,
translations, or follow-ups of each other.  The \verb|\relatedentry|
command can render the relationship.

The second edition of Knuth's fundamental work
\cite{knuthFundamentalAlgorithms1973} is related to the original TeXbook
\cite{knuthTeXbook1986}---both are part of the same author's
monumental series on computer programming.

% =====================================================
\section{Sorting Options}
\label{sec:sorting}
% =====================================================

The bibliography can be sorted in different orders.  The current
configuration uses \texttt{sorting=nyt} (name, year, title), which is the
standard for author-year styles.  Alternative sortings include:

\begin{description}
    \item[nyt] Name, Year, Title --- default for authoryear.
    \item[ynt] Year, Name, Title --- chronological ordering.
    \item[nyvt] Name, Year, Volume, Title --- for series.
    \item[none] Citation order --- no sorting, entries appear in the order
        they are first cited.
\end{description}

To switch, change the \texttt{sorting} option in the \verb|\usepackage|
declaration and recompile.

% =====================================================
\section{BibLaTeX Options in Detail}
\label{sec:options}
% =====================================================

\subsection{Backend Selection}

BibLaTeX supports two backends: \texttt{biber} (modern, recommended) and
\texttt{bibtex} (legacy).  Biber provides full Unicode support, advanced
sorting, and the \texttt{related} field.  The current configuration uses
biber.

\subsection{Dashed Repeated Authors}

With \texttt{dashed=true}, repeated author names in the bibliography are
replaced by a dash (\textemdash), reducing visual clutter.  See the full
bibliography below for examples where Knuth's multiple works appear with
dashes.

\subsection{Name Formatting}

BibLaTeX offers fine-grained control over name formatting.  The
\texttt{maxbibnames} and \texttt{minbibnames} options control when
et~al.~truncation applies.  The current setting shows all names
(\texttt{maxbibnames=99}).

\subsection{Field Inclusion}

The \texttt{doi}, \texttt{isbn}, \texttt{url}, and \texttt{eprint}
options control which identifier fields appear in the bibliography.
Disabling \texttt{isbn} and \texttt{url} produces cleaner entries when
those fields are not essential.

% =====================================================
\section{Compatibility and Multi-Source}
\label{sec:multi-source}
% =====================================================

BibLaTeX can coexist with natbib-style commands through the
\texttt{natbib} compatibility option.  While not enabled in this
configuration, adding \texttt{natbib} to the package options provides
access to commands like \verb|\citet| and \verb|\citep|.

Multiple \verb|\addbibresource| calls allow combining bibliographies from
different sources.  This document loads both the main repository
bibliography (\texttt{../../bib/bibliography.bib}) and a local
supplementary file (\texttt{local-refs.bib}) containing custom entry types.

% =====================================================
\section{Bibliography Filters}
\label{sec:filters}
% =====================================================

Beyond type and keyword filtering, BibLaTeX supports exclusion filters.
Print everything \emph{except} articles:
\printbibliography[nottype=article,heading=subbibliography,title={Non-Article Entries}]

Exclude entries by name or keyword:
\printbibliography[notkeyword=reference,heading=subbibliography,title={Non-Reference Entries}]

% =====================================================
\section{Bibliography Headings}
\label{sec:headings}
% =====================================================

BibLaTeX provides several heading styles for bibliography sections:

\begin{description}
    \item[default] The top-level heading, usually ``References'' or
        ``Bibliography''.
    \item[subbibliography] A subordinate heading that does not appear in
        the table of contents.
    \item[subbibintoc] A subordinate heading that \emph{does} appear in
        the table of contents.
    \item[bibliography] Explicit bibliography heading, useful when
        overriding the default.
\end{description}

The choice of heading affects how deeply the bibliography is nested in the
document structure.  For articles and reports, \texttt{subbibliography} is
typical.  For books and theses, \texttt{subbibintoc} provides better
navigation.

% =====================================================
\section{Name List Customisation}
\label{sec:name-lists}
% =====================================================

BibLaTeX offers precise control over how author and editor lists are
displayed.  The current configuration sets \texttt{maxbibnames=99} so
all names appear in the bibliography.  With a lower value, such as 3,
lists exceeding that length would be truncated to the first name
plus ``et~al.''

The \texttt{uniquelist} option, when enabled, adds extra names from the
list to disambiguate entries by the same author in the same year.
Disabling it (\texttt{uniquelist=false}) keeps names consistent across
entries.

The \texttt{minbibnames} threshold determines the minimum number of names
shown before truncation occurs.  Setting it equal to \texttt{maxbibnames}
ensures no truncation happens regardless of list length.

% =====================================================
\section{Shorthand Entries}
\label{sec:shorthand}
% =====================================================

Some references have well-known short forms.  The \texttt{shorthand}
field defines an abbreviated citation label.  For example, the PTB SI
units reference uses the shorthand \texttt{PTB07}, and the GBV catalog
is cited as \texttt{GBV}.

After the first citation, subsequent references use the shorthand:
\cite{PTBSI}.

The shorthand is particularly useful for frequently cited standards,
handbooks, and institutional references where the full citation would be
cumbersome in running text.

% =====================================================
\section{Abstracts and Annotations}
\label{sec:abstracts}
% =====================================================

BibLaTeX can display the \texttt{abstract} and \texttt{annotation} fields
from bibliography entries.  While not rendered in the standard bibliography
output, these fields are accessible through custom bibliography drivers or
verbose citation commands.

The McKinney reference includes an abstract describing Python data analysis
techniques\footnote{See the full abstract in the bibliography entry for
\citetitle{mckinneyPythonDataAnalysis2018}.}.  Similarly, the Ramalho
reference provides a detailed abstract about idiomatic Python
programming\footnote{The abstract for \citetitle{ramalhoFluentPython2015}
describes leveraging Python's best features.}.

% =====================================================
\section{Language and Localization}
\label{sec:localization}
% =====================================================

BibLaTeX respects the document's language setting for bibliography strings
such as ``and'', ``editor'', ``In'', and ``volume''.  With
\texttt{language=english}, all biblatex strings appear in English.

The \texttt{langid} field in individual entries overrides the document
language for that entry.  The Baehr thermodynamics reference
\cite{baehrThermodynamikGrundlagenUnd2016} has \texttt{langid=german},
which can be used for language-specific sorting or formatting.

BibLaTeX automatically translates strings based on the active language.
Switching to German in the document preamble would produce ``Herausgeber''
instead of ``editor'', ``Band'' instead of ``volume'', and so on.

% =====================================================
\section{DOIs and Identifier Links}
\label{sec:identifiers}
% =====================================================

The \texttt{doi} option controls whether Digital Object Identifiers are
rendered as clickable links.  With \texttt{doi=true}, the Einstein paper
shows its DOI: \cite{einstein1905}.  The DOI field is automatically
formatted as a URL pointing to \texttt{https://doi.org/...}.

The \texttt{isbn} and \texttt{url} options are disabled in this
configuration to keep the bibliography clean.  Re-enabling them would add
ISBN numbers for books and URLs for online entries.

The \texttt{eprint} option controls the display of arXiv identifiers.
When enabled, preprint references show their arXiv ID with a link to the
arXiv abstract page.

% =====================================================
\section{Cross-Reference Examples}
\label{sec:crossref}
% =====================================================

OmniLaTeX integrates \texttt{cleveref} for intelligent cross-references:

\begin{itemize}
    \item \verb|\cref{sec:citation-styles}| $\to$ \cref{sec:citation-styles}
    \item \verb|\cref{sec:qualifiers}| $\to$ \cref{sec:qualifiers}
    \item \verb|\cref{sec:multiple}| $\to$ \cref{sec:multiple}
    \item \verb|\cref{sec:date-tracking}| $\to$ \cref{sec:date-tracking}
    \item \verb|\cref{sec:citation-counting}| $\to$ \cref{sec:citation-counting}
    \item \verb|\cref{sec:categories}| $\to$ \cref{sec:categories}
    \item \verb|\cref{sec:inline}| $\to$ \cref{sec:inline}
    \item \verb|\cref{sec:custom-types}| $\to$ \cref{sec:custom-types}
    \item \verb|\cref{sec:related}| $\to$ \cref{sec:related}
    \item \verb|\cref{sec:sorting}| $\to$ \cref{sec:sorting}
    \item \verb|\cref{sec:options}| $\to$ \cref{sec:options}
    \item \verb|\cref{sec:filters}| $\to$ \cref{sec:filters}
    \item \verb|\cref{sec:name-lists}| $\to$ \cref{sec:name-lists}
    \item \verb|\cref{sec:shorthand}| $\to$ \cref{sec:shorthand}
    \item \verb|\cref{sec:localization}| $\to$ \cref{sec:localization}
    \item \verb|\cref{sec:identifiers}| $\to$ \cref{sec:identifiers}
\end{itemize}

% =====================================================
% Bibliography
% =====================================================
\printbibliography

\end{document}
